\documentclass{sprz}

\begin{document}

\tableofcontents

\chapter{Raport z laboratorium CI}

\section{Cel}

Celem laboratorium było zbudowanie potoku ciągłej integracji dla projektu
Mission Control, który składa się z backendu w~.NET oraz frontendu w~Angularze.
Potok ma szybko informować o~regresjach, nie marnować czasu runnerów na
niepotrzebne buildy i~jednoznacznie odpowiadać na pytanie, czy zmianę można
scalić do gałęzi \texttt{main}.

\section{Detekcja zmienionych komponentów}

Job \texttt{changes} pobiera przez API GitHuba pełną, stronicowaną listę plików
zmienionych w~pull requeście i~klasyfikuje je jako zmiany backendu, frontendu
albo dokumentacji. Zmiany w~katalogu \texttt{.github} konserwatywnie wymagają
sprawdzenia obu komponentów, a~nieznana ścieżka nigdy nie daje dwóch wartości
\texttt{false}. Dla zdarzeń \texttt{push} i~ręcznego startu oba komponenty są
zawsze sprawdzane.

\section{Jeden zawsze obecny gate}

Komponenty uruchamiają się warunkowo, dlatego ochrona gałęzi \texttt{main} nie
wymaga już ich checków bezpośrednio. Wymagany jest tylko job
\texttt{ci-required}, który ocenia wynik detekcji oraz obu komponentów.
Komponent wymagany musi zakończyć się sukcesem, a~niewymagany musi zostać
pominięty. Każdy inny wynik, także błąd samej detekcji, daje czerwony gate
i~blokuje scalenie.

\section{Reużywanie i~osobne uruchomienia}

Procedura backendu została przeniesiona do akcji typu composite, a~cały job
backendu do reusable workflow z~typowanymi wejściami. Osobny workflow
\texttt{Inspect CI} uruchamiany zdarzeniem \texttt{workflow\_run} pobiera
artefakt metadanych z~dokładnie wskazanego runu producenta i~sprawdza jego
schemat oraz zgodność SHA, identyfikatora runu i~numeru próby.

\section{Wydawanie dokumentu PDF}

Ten raport jest kompilowany w~CI przy każdym pull requeście i~każdym pushu na
gałąź \texttt{main}. Kompilacja ma wyłącznie prawo odczytu, a~klasa szablonu
jest pobierana z~przypiętej rewizji publicznego repozytorium. Pull request
zmieniający tylko pliki TeX publikuje wersję wstępną (prerelease), a~push na
\texttt{main} pełne wydanie. Pull request zmieniający definicje workflow albo
pochodzący z~forka kończy się samym artefaktem, bez zapisu wydania.

Nazwa każdego wydania zawiera identyfikator runu i~numer próby, dlatego
ponowne uruchomienie tworzy nowy tag zamiast nadpisywać poprzedni. Tag zawsze
wskazuje dokładnie ten commit, z~którego skompilowano dokument.

\section{Wyniki}

\begin{tabular}{|l|l|}
\hline
\textbf{Scenariusz} & \textbf{Wynik gate'a} \\
\hline
Zmiana tylko dokumentacji & zielony, oba komponenty pominięte \\
Zmiana tylko backendu & zielony, frontend pominięty \\
Zmiana workflow & zielony, oba komponenty uruchomione \\
Regresja w~backendzie & czerwony, scalenie zablokowane \\
Uszkodzona detekcja & czerwony, scalenie zablokowane \\
\hline
\end{tabular}

\section{Wnioski}

Optymalizacja zakresu uruchamianych jobów nie może zmieniać znaczenia
komunikatu ,,CI przeszło''. Jeden jednoznaczny gate oddziela politykę scalania
od liczby faktycznie uruchomionych jobów, a~kontrolowane błędy na osobnych
gałęziach pozwalają sprawdzić, że każda czerwona ścieżka rzeczywiście blokuje
scalenie.

TEST
\end{document}
